\section{Task 2: KDEF dataset study}
\label{sec:task2}

In this task we take the best FER2013 representation from Task~1, HOG, retrain the three classifiers on KDEF, and compare the results with FER2013.

The Karolinska Directed Emotional Faces (KDEF) dataset contains 4,900 pictures of 70 individuals displaying 7 emotions, photographed in a controlled studio setup. The pose of most interest to us is \textit{straight}; the full collection protocol is described in Appendix~\ref{app:kdef-desc}.

KDEF images ($562\times762$ RGB) are converted to the FER2013 format ($48\times48$ grayscale, scaled to $[0,1]$) with the common representation of Section~\ref{sec:data}; the face-detection details are given in Appendix~\ref{app:kdef-desc}.

The data is split 80/20 by subject i.e out of the 70 subjects, 56 of them are used in the training set, and 14 are in the test set.  This is to ensure there is no data leakage into the test set. KDEF is perfectly balanced, with 112 training and 28 test images per emotion (Figure~\ref{fig:kdef-split} in Appendix~\ref{app:figs}).

In the previous task, we found out that the best feature representation is Histogram of Oriented Gradients (HOG), so every KDEF image is described with the same 900 HOG features (Section~\ref{sec:data}).


\subsection{Model training and selection}

We train the same three classifiers as in Task 1: Logistic Regression, SVM and Random Forest. For each classifier we run a grid search with 5-fold cross-validation, where the folds are also split by subject (using GroupKFold), so the model is always validated on subjects it has not seen during training. The search space (20, 96 and 84 configurations) is listed in Table~\ref{tab:kdef-grid} in Appendix~\ref{app:tables}; the selected settings are shown in bold in that table.


\subsection{Results}

\begin{table}[h]\centering\small
\begin{tabular}{@{}lcccc@{}}\toprule
 & \multicolumn{2}{c}{FER2013} & \multicolumn{2}{c}{KDEF} \\\cmidrule(lr){2-3}\cmidrule(l){4-5}
Classifier & Validation & Test & Validation & Test \\\midrule
Logistic Regression & 0.395 & 0.378 & 0.669 & 0.581 \\
SVM & \textbf{0.505} & \textbf{0.568} & \textbf{0.679} & 0.604 \\
Random Forest & 0.440 & -- & 0.611 & \textbf{0.614} \\
\bottomrule\end{tabular}
\caption{Macro-F1 with HOG features, for the models trained and evaluated on FER2013 (Task~1) and on KDEF.}
\label{tab:kdef}

\end{table}
Table~\ref{tab:kdef} shows that all three classifiers score higher on KDEF than on FER2013, even though they are trained on only 784 images instead of about 28{,}700. The reason is that KDEF is a much cleaner dataset. Every photo is taken in the same studio, with the same lighting and camera position, and the actors show each emotion clearly, so the faces line up well and the expression is almost the only thing that differs between images. FER2013 images come from the web, with different lighting, head poses, faces partly covered by hands or glasses, and labels that are sometimes wrong, so the classifiers have to learn the same emotion from very different-looking images.


\subsection{Assignment Questions}
\textbf{Why are KDEF images easier or harder to classify?} Overall, KDEF is easier: with the same HOG features, the SVM reaches a macro-F1 of 0.679 in cross-validation and 0.604 on the test subjects. The consistency of the KDEF images that makes this possible is described in Appendix~\ref{app:kdef-consistency}.

The mean face per emotion for both datasets (Figure~\ref{fig:mean-faces}, Appendix~\ref{app:figs}) shows how consistently each facial feature is positioned across images: the blurrier the mean face, the less consistent the images for that emotion.

\textbf{Which models change ranking?} Table~\ref{tab:kdef-model-ranking-comparison} compares the rankings.

\begin{table}[h]\centering\small
\begin{tabular}{@{}clll@{}}\toprule
Rank & FER2013 validation & KDEF cross-validation & KDEF test \\\midrule
1 & SVM (0.505) & SVM (0.679) & Random Forest (0.614) \\
2 & Random Forest (0.440) & Logistic Regression (0.669) & SVM (0.604) \\
3 & Logistic Regression (0.395) & Random Forest (0.611) & Logistic Regression (0.581) \\
\bottomrule\end{tabular}
\caption{FER2013 vs KDEF model comparison}
\label{tab:kdef-model-ranking-comparison}
\end{table}

On KDEF, Logistic Regression jumps from last to second place in validation, while the Random Forest falls to last. We think this is because the well-aligned KDEF faces make the HOG features almost linearly separable, so a strongly regularized linear model ($C=0.01$) works well with few training images. Although, the Random Forest performs better in the test set by a small margin, we choose the best model from the validation set.


\textbf{Which emotions transfer well between datasets?} To see which emotions hold up on both datasets, we compare the F1 score per emotion of the KDEF models with that of the FER-trained SVM from Task~1 (Table~\ref{tab:emotion-both} in Appendix~\ref{app:tables}). Happy is the easiest emotion on both datasets (0.93 on KDEF and 0.75 on FER2013 for the SVM), and Surprise also scores well on both (0.62 and 0.73). Both expressions have large, distinctive shapes, a smile or an open mouth with raised eyebrows, which HOG picks up whether the expression is posed in a studio or caught in a web photo. Disgust scores well on both too, but FER2013 contains very few Disgust images (1.5\%), so this result is less certain. Fear, Sad and Angry are among the hardest emotions on both datasets, with F1 scores between 0.33 and 0.55.

The confusion matrices in Appendix~\ref{app:cm} (Figure~\ref{fig:kdef-confusion}) show that most mistakes happen between emotions that look alike. Fear is the clearest example: the SVM mistakes more fearful faces for Surprise (43\%) than it gets right (36\%), and the other two classifiers do the same. This is not surprising, since both expressions have wide-open eyes, raised eyebrows and an open mouth. In the same way, Sad is often mistaken for Neutral (21\%), and Angry for Fear (25\%) or Disgust (14\%). Based on this, we expect Happy and Surprise to transfer best between the datasets, and Fear, Sad and Angry to transfer worst.
